\documentclass[12pt, a4paper]{article}

% Gold-Standard Academic Typography & Encoding Packages
\usepackage[T1]{fontenc}
\usepackage{lmodern}
\usepackage{microtype}
\usepackage[utf8]{inputenc}
\usepackage[margin=1in]{geometry}
\setlength{\headheight}{18pt}
\usepackage{graphicx}
\usepackage{hyperref}
\usepackage{amsmath, amssymb}
\usepackage{booktabs}
\usepackage{float}
\usepackage[export]{adjustbox}
\usepackage{listings}
\usepackage{enumitem}
\usepackage{xcolor}
\usepackage{fancyhdr}
\usepackage{array}
\usepackage{tabularx}
\usepackage{longtable}
\usepackage{caption}
\usepackage{subcaption}

% Hyperref Setup (Print-Safe & Digital PDF Friendly)
\hypersetup{
    colorlinks=true,
    linkcolor=black,
    citecolor=black,
    urlcolor=black,
    pdftitle={SRS - CampusCare (CSE 416)},
    pdfauthor={Rudro Antony Mrong, Md. Masud Rahman}
}

% Header and Footer Setup (No Overlapping Text)
\pagestyle{fancy}
\fancyhf{}
\rhead{\small \textbf{CSE 416: Software Project Design}}
\lhead{\small \textbf{CampusCare SRS}}
\rfoot{\small Page \thepage}
\lfoot{\small UITS - Dept. of CSE}
\renewcommand{\headrulewidth}{0.4pt}
\renewcommand{\footrulewidth}{0.4pt}

% Code Listing Style Setup
\lstset{
    basicstyle=\ttfamily\small,
    breaklines=true,
    frame=single,
    backgroundcolor=\color{gray!5},
    keywordstyle=\color{blue},
    commentstyle=\color{green!60!black},
    stringstyle=\color{purple}
}

\begin{document}

% =========================================================
% COVER PAGE
% =========================================================
\begin{titlepage}
    \centering
    
    % University Logo
    
\includegraphics[width=0.28\textwidth]{uits_logo.png}\\[1.2cm]
    
    {\Large \textbf{University of Information Technology and Sciences (UITS)}}\\[0.3cm]
    {\large Department of Computer Science and Engineering}\\[1.2cm]
    
    {\Large \textbf{COURSE:}} Software Project Design and Development (CSE 416)\\[1.0cm]
    
    \rule{\linewidth}{0.5mm}\\[0.4cm]
    {\LARGE \textbf{SOFTWARE REQUIREMENTS SPECIFICATION (SRS)}}\\[0.2cm]
    {\large \textbf{Project Title: CampusCare}}\\[0.1cm]
    {\small Centralized University Complaint Management System}\\[0.4cm]
    \rule{\linewidth}{0.5mm}\\[1.5cm]
    
    \begin{minipage}{0.48\textwidth}
        \begin{flushleft} \large
            \textbf{Submitted To:}\\[0.2cm]
            \textbf{Al-Imtiaz}\\
            Associate Professor \& Head\\
            Ph.D. (Research Fellow), BUET\\
            Department of CSE, UITS
        \end{flushleft}
    \end{minipage}
    \hfill
    \begin{minipage}{0.48\textwidth}
        \begin{flushright} \large
            \textbf{Submitted By (Group Members):}\\[0.2cm]
            \textbf{Rudro Antony Mrong}\\
            ID: 0432320005101059\\[0.2cm]
            \textbf{Md. Masud Rahman}\\
            ID: 0432320005101064
        \end{flushright}
    \end{minipage}
    
    \vfill
    {\large Submission Date: August 2026}
\end{titlepage}

% =========================================================
% ABSTRACT PAGE
% =========================================================
\newpage
\pagenumbering{roman}
\setcounter{page}{2}

\begin{abstract}
\noindent
CampusCare is an internal, web-based 3-tier ticketing and facility helpdesk application designed specifically for university campus infrastructure management. In many higher education institutions, maintenance reporting relies on unorganized communication channels such as verbal walk-ins, phone calls, and unmonitored emails. This fragmentation leads to misplaced requests, delayed issue resolutions, zero progress visibility for reporters, and a complete lack of administrative oversight regarding institutional asset health. 

CampusCare addresses these operational challenges by implementing a centralized digital pipeline governed by a strict state machine workflow (\texttt{Submitted} $\rightarrow$ \texttt{Assigned} $\rightarrow$ \texttt{In Progress} $\rightarrow$ \texttt{Resolved} $\rightarrow$ \texttt{Closed} / \texttt{Reopened}). Built using modern full-stack web technologies—including Next.js (App Router, React 19, TypeScript), Tailwind CSS, Shadcn UI, and Supabase (PostgreSQL engine with Row Level Security policies and Supabase Storage)—the platform provides isolated, role-based interfaces for Students (End-Users), Staff (Resolvers), and Administrators (Dispatchers \& Analysts).

This Software Requirements Specification (SRS) document presents a comprehensive analysis of CampusCare, detailing the functional requirements, non-functional performance benchmarks, relational database entity schemas, behavioral system models (Use Case, Class, State Machine, ER, Sequence, and Data Flow Diagrams), and a traceably mapped system testing plan.
\end{abstract}

\newpage
\tableofcontents
\newpage
\pagenumbering{arabic}

% =========================================================
% SECTION 1: INTRODUCTION
% =========================================================
\section{Introduction}

\subsection{Purpose}
The purpose of this Software Requirements Specification (SRS) document is to define the detailed functional, non-functional, interface, database, and behavioral requirements for \textbf{CampusCare: Centralized University Complaint Management System}. This document follows the IEEE 830 standard guidelines and acts as an authoritative baseline for software engineers, database architects, UI/UX designers, quality assurance testers, and project instructors.

\subsection{Document Conventions}
This document observes standard typography and formatting conventions:
\begin{itemize}[leftmargin=*]
    \item \textbf{Bold Text:} Indicates UI labels, role designations, entity names, or critical emphasis.
    \item \texttt{Monospace Text:} Denotes database table names, SQL attributes, route paths, environment parameters, or state machine statuses.
    \item \textbf{Requirement Identifiers:} All functional and non-functional requirements observe structured, readable numbering schemes (e.g., \textbf{Requirement 4.1.A}, \textbf{Requirement 5.3.B}) to facilitate backward and forward traceability in testing.
\end{itemize}

\subsection{Definitions, Acronyms, and Abbreviations}
To facilitate clarity across technical and evaluation audiences, the acronyms and terms used throughout this document are defined in Table \ref{tab:acronyms}.

\begin{table}[H]
    \centering
    \caption{Centralized Glossary of Definitions, Acronyms, and Abbreviations}
    \label{tab:acronyms}
    \begin{tabularx}{\textwidth}{l X}
        \toprule
        \textbf{Term / Acronym} & \textbf{Definition} \\
        \midrule
        \textbf{SRS} & Software Requirements Specification (IEEE 830 standard document). \\
        \textbf{RBAC} & Role-Based Access Control (Restricting access based on user role). \\
        \textbf{SLA} & Service Level Agreement (Target resolution timeframe based on urgency). \\
        \textbf{RLS} & Row Level Security (PostgreSQL engine-level database access control). \\
        \textbf{RSC} & React Server Components (Next.js server-rendered component paradigm). \\
        \textbf{UUID} & Universally Unique Identifier (128-bit database primary key standard). \\
        \textbf{JSONB} & Binary JSON storage format in PostgreSQL supporting indexed sub-queries. \\
        \textbf{KPI} & Key Performance Indicator (Aggregate quantitative performance metric). \\
        \textbf{DFD} & Data Flow Diagram (Process and data store architecture model). \\
        \textbf{ERD} & Entity-Relationship Diagram (Database schema relationship model). \\
        \textbf{UITS} & University of Information Technology and Sciences. \\
        \textbf{BUET} & Bangladesh University of Engineering and Technology. \\
        \bottomrule
    \end{tabularx}
\end{table}

\subsection{Intended Audience and Reading Suggestions}
This document is designed for the following stakeholders:
\begin{itemize}[leftmargin=*]
    \item \textbf{Course Instructor \& Evaluators:} To verify project scope, engineering rigor, requirement completeness, and design correctness for CSE 416.
    \item \textbf{Frontend \& Backend Developers:} To guide UI component implementation, Next.js Server Actions, and Supabase database interactions.
    \item \textbf{Quality Assurance Testers:} To execute test plans against the traceability matrix in Section 8.
\end{itemize}

\subsection{Product Scope}
CampusCare is a centralized university facility helpdesk system designed to replace manual, paper-based, or fragmented communication channels. The product scope includes:
\begin{enumerate}[leftmargin=*]
    \item \textbf{Student Module:} Authenticated complaint submission with optional image upload, human-readable ID generation (\texttt{CMP-YYYY-XXXX}), real-time status timeline tracking, and resolution confirmation/rejection loop with optional anonymous star ratings.
    \item \textbf{Staff Module:} Departmental task queue, SLA deadline indicators, status transition controls, repair completion proof upload, and progress note feeds.
    \item \textbf{Admin Module:} Global command center with multi-parameter ticket search/filtering, manual ticket assignment, priority escalation, user management, category/department taxonomies, analytics charts, and immutable audit logs.
    \item \textbf{Data Core:} Relational PostgreSQL database with Row Level Security (RLS), foreign key integrity, and trigger-driven audit logging.
\end{enumerate}

\subsection{References}
\begin{enumerate}[leftmargin=*]
    \item IEEE Std 830-1998, \textit{IEEE Recommended Practice for Software Requirements Specifications}.
    \item Next.js App Router Documentation: \url{https://nextjs.org/docs}
    \item Supabase PostgreSQL \& RLS Security Guide: \url{https://supabase.com/docs/guides/auth/row-level-security}
    \item CampusCare Product Requirements Document (\texttt{doc/PRD.md}).
    \item CampusCare Project Proposal (\texttt{doc/Project\_Proposal.md}).
\end{enumerate}

% =========================================================
% SECTION 2: OVERALL DESCRIPTION
% =========================================================
\section{Overall Description}

\subsection{Product Perspective}
CampusCare is a standalone, web-based application built using a 3-tier architecture:
\begin{itemize}[leftmargin=*]
    \item \textbf{Presentation Tier:} Next.js (App Router, React 19, TypeScript) rendered with Tailwind CSS and Shadcn UI components.
    \item \textbf{Application Tier:} Server-side business logic, API endpoints, and Next.js Server Actions enforcing Zod input validation schemas.
    \item \textbf{Data Tier:} PostgreSQL database hosted via Supabase enforcing engine-level Row Level Security policies and Supabase Storage buckets for photographic attachments.
\end{itemize}

\subsection{Product Functions}
The principal functions of CampusCare are summarized below:
\begin{itemize}[leftmargin=*]
    \item User Authentication \& Session Management via \texttt{HttpOnly} cookies.
    \item Categorized Complaint Lodging with optional photo evidence attachment.
    \item Unique Complaint ID auto-generation (\texttt{CMP-YYYY-XXXX}).
    \item Departmental Ticket Queue Routing \& SLA Deadline Calculation.
    \item Interactive Status Progression (\texttt{Submitted} $\rightarrow$ \texttt{Assigned} $\rightarrow$ \texttt{In Progress} $\rightarrow$ \texttt{Resolved} $\rightarrow$ \texttt{Closed} / \texttt{Reopened}).
    \item Resolution Verification Feedback Loop (Student Confirmation \& Optional Anonymous 1--5 Star Rating).
    \item Administrative Master Data Dispatching, Priority Override, and Taxonomy Configuration.
    \item Immutable PostgreSQL Database Trigger Audit Logging.
\end{itemize}

\subsection{User Classes and Characteristics}
The system isolates operational permissions into three primary user classes, summarized with wrapped text in Table \ref{tab:user_classes}.

\begin{table}[H]
    \centering
    \caption{User Class Characteristics and System Permissions}
    \label{tab:user_classes}
    \begin{tabularx}{\textwidth}{p{0.18\textwidth} p{0.25\textwidth} X}
        \toprule
        \textbf{User Class} & \textbf{Target Persona} & \textbf{System Capabilities \& Access Boundaries} \\
        \midrule
        \textbf{Student} & Campus Students & Submit infrastructure issues, upload optional photos, track timelines, confirm or reject resolution actions, and provide optional anonymous feedback. Access restricted to personal records via RLS. \\
        \addlinespace
        \textbf{Staff} & Maintenance \& IT Personnel & View departmental task queue, accept assigned tickets, append chronological progress notes, upload repair completion proof, and mark tickets as resolved. Access restricted to departmental tickets. \\
        \addlinespace
        \textbf{Admin} & Facility Managers \& Dept Heads & Dispatch unassigned tickets, reassign personnel, override SLA priority levels, manage user accounts/taxonomies, analyze system KPI metrics, and review read-only audit logs. Full system read/write access. \\
        \bottomrule
    \end{tabularx}
\end{table}

\subsection{Operating Environment}
\begin{itemize}[leftmargin=*]
    \item \textbf{Client Side:} Modern desktop and mobile web browsers (Chrome, Safari, Firefox, Edge) with JavaScript enabled. Minimum viewport width: 320px.
    \item \textbf{Server Side:} Node.js 20+ execution environment hosting Next.js.
    \item \textbf{Database Engine:} PostgreSQL 15+ hosted via Supabase Cloud.
\end{itemize}

\subsection{Design and Implementation Constraints}
\begin{itemize}[leftmargin=*]
    \item \textbf{Security Constraint:} Authentication tokens must never be exposed to client JavaScript or stored in \texttt{LocalStorage}. Tokens are stored strictly in \texttt{HttpOnly}, \texttt{SameSite=Lax}, \texttt{Secure} cookies.
    \item \textbf{Data Isolation Constraint:} PostgreSQL Row Level Security (RLS) policies must be enabled on 100\% of database tables to prevent unauthorized cross-tenant data access.
    \item \textbf{File Upload Constraint:} Image attachments are strictly limited to JPEG, PNG, and WEBP formats with a maximum file size of 5MB per upload.
\end{itemize}

\subsection{User Documentation}
User documentation will include online dashboard tooltips, inline form field guides, and a digital PDF User Manual for Staff Resolvers and Administrators.

\subsection{Assumptions and Dependencies}
\begin{itemize}[leftmargin=*]
    \item \textbf{University Email Integration:} The system assumes students and staff register and log in using their official university email addresses.
    \item \textbf{Cloud Service Availability:} Supabase Cloud infrastructure (PostgreSQL database, Storage buckets, and Authentication endpoints) maintains an operational availability exceeding 99.5\% uptime.
    \item \textbf{Network Infrastructure:} Campus Wi-Fi and mobile networks are available across university premises.
\end{itemize}

% =========================================================
% SECTION 3: EXTERNAL INTERFACE REQUIREMENTS
% =========================================================
\section{External Interface Requirements}

\subsection{User Interfaces}
The application exposes three responsive user interface views:
\begin{enumerate}[leftmargin=*]
    \item \textbf{Student Dashboard View (\texttt{/student/dashboard}):} Features metric cards (Total, Active, Resolved), submission modal button, and active ticket list with status timeline drawers.
    \item \textbf{Staff Workspace View (\texttt{/staff/dashboard}):} Features departmental queue cards, SLA deadline countdown timers, progress note textareas, and completion photo dropzone.
    \item \textbf{Admin Command Center (\texttt{/admin/dashboard}):} Features system KPI cards, multi-parameter search/filter master table, dispatcher modal overlays, category manager, and read-only audit log tables.
\end{enumerate}

\subsection{Hardware Interfaces}
No specialized hardware is required. The web application interfaces with standard desktop, tablet, and mobile device displays and webcams/cameras for photo capturing.

\subsection{Software Interfaces}
\begin{itemize}[leftmargin=*]
    \item \textbf{Supabase Auth Service:} Interface for credentials validation, session JWT issuance, and cookie synchronization.
    \item \textbf{Supabase Storage API:} Interface for secure binary file uploads and signed URL retrieval.
\end{itemize}

\subsection{Communications Interfaces}
The platform enforces encrypted Transport Layer Security (TLS 1.3) across all client-server communications via HTTPS. API payloads observe standard JSON format over HTTP POST/GET requests.

% =========================================================
% SECTION 4: SYSTEM FEATURES
% =========================================================
\section{System Features (Functional Subsystems)}

\subsection{Student Module: Complaint Submission \& Tracking Engine}

\subsubsection{Description and Priority}
This subsystem allows authenticated students to lodge facility complaints with detailed location descriptions and optional photo attachments, track status changes in real-time, and view progress notes. Priority: \textbf{High}.

\subsubsection{Stimulus / Response Sequences}
\begin{itemize}[leftmargin=*]
    \item \textbf{Stimulus:} Student clicks "Submit Complaint" and fills out form fields.
    \item \textbf{Response:} System validates payload via Zod, uploads optional evidence to storage, generates a unique \texttt{CMP-YYYY-XXXX} ID, inserts database record, and displays success receipt.
    \item \textbf{Stimulus:} Student selects a ticket from their dashboard list.
    \item \textbf{Response:} System fetches ticket record and renders a chronological tracking timeline with progress notes and assigned staff details.
\end{itemize}

\subsubsection{Functional Requirements}
\begin{itemize}[leftmargin=*]
    \item \textbf{Requirement 4.1.A - Session Authentication:} The system shall restrict student portal access to authenticated users possessing the \texttt{student} role.
    \item \textbf{Requirement 4.1.B - Complaint Form Fields:} The system shall require title, category selection, location details, priority level, and detailed description during submission.
    \item \textbf{Requirement 4.1.C - Photo Upload Validation (Optional):} The system shall support optional image attachments, validating files for type (JPG, PNG, WEBP) and maximum size (5MB).
    \item \textbf{Requirement 4.1.D - Ticket Number Generation:} The system shall automatically execute a database trigger function generating a human-readable Complaint ID (e.g., \texttt{CMP-2026-0001}).
    \item \textbf{Requirement 4.1.E - Personal Data Isolation:} The system shall prevent students from viewing or searching complaints lodged by other students via PostgreSQL RLS policies.
    \item \textbf{Requirement 4.1.F - Real-Time Timeline Rendering:} The system shall render a chronological timeline displaying ticket creation, staff assignment, progress updates, and resolution status.
\end{itemize}

\subsection{Staff Module: Task Queue \& SLA Resolution Workflow}

\subsubsection{Description and Priority}
This subsystem provides maintenance and technical staff with a queue of tickets assigned to their department, SLA warning indicators, progress note submission forms, repair proof uploads, and status transition toggles. Priority: \textbf{High}.

\subsubsection{Stimulus / Response Sequences}
\begin{itemize}[leftmargin=*]
    \item \textbf{Stimulus:} Staff member logs in and opens the departmental dashboard.
    \item \textbf{Response:} System queries database for complaints matching the staff member's department ID and renders a prioritized queue with SLA badges.
    \item \textbf{Stimulus:} Staff member clicks "Accept Ticket".
    \item \textbf{Response:} System updates ticket status from \texttt{Assigned} to \texttt{In Progress} and records an audit log entry.
    \item \textbf{Stimulus:} Staff member appends a note and uploads a repair photo, then clicks "Mark Resolved".
    \item \textbf{Response:} System updates ticket status to \texttt{Resolved}, attaches repair proof, and notifies the reporting student.
\end{itemize}

\subsubsection{Functional Requirements}
\begin{itemize}[leftmargin=*]
    \item \textbf{Requirement 4.2.A - Department Queue Filtering:} The system shall restrict staff queue views strictly to tickets matching their assigned operational department.
    \item \textbf{Requirement 4.2.B - SLA Target Display:} The system shall display visual warning indicators for tickets approaching or exceeding target SLA resolution deadlines.
    \item \textbf{Requirement 4.2.C - State Transition to In Progress:} The system shall allow staff members to accept assigned tickets, transitioning state to \texttt{In Progress}.
    \item \textbf{Requirement 4.2.D - Progress Note Logging:} The system shall allow staff to append public or internal chronological notes to active tickets.
    \item \textbf{Requirement 4.2.E - Repair Proof Upload \& Resolution:} The system shall support photographic repair completion uploads and allow staff to mark tickets as \texttt{Resolved}.
\end{itemize}

\subsection{Admin Module: Dispatcher Command Center \& Audit Ledger}

\subsubsection{Description and Priority}
This subsystem provides facility managers with global ticket routing, staff re-assignment, priority escalation, user management, category administration, analytics charts, and immutable audit logs. Priority: \textbf{High}.

\subsubsection{Stimulus / Response Sequences}
\begin{itemize}[leftmargin=*]
    \item \textbf{Stimulus:} Admin opens master complaint table and selects an unassigned ticket.
    \item \textbf{Response:} System opens routing modal displaying active staff members and departments.
    \item \textbf{Stimulus:} Admin assigns staff member and sets priority to \texttt{Critical}.
    \item \textbf{Response:} System updates complaint record, recalculates SLA target time, and writes an audit log.
\end{itemize}

\subsubsection{Functional Requirements}
\begin{itemize}[leftmargin=*]
    \item \textbf{Requirement 4.3.A - Master Data Table Filtering:} The system shall provide a master table supporting text search, status/priority/department filtering, sorting, and pagination.
    \item \textbf{Requirement 4.3.B - Staff Ticket Dispatching:} The system shall allow administrators to assign unassigned tickets or reassign active tickets across departments and personnel.
    \item \textbf{Requirement 4.3.C - Priority Override:} The system shall allow administrators to override and elevate or lower the priority level of any complaint.
    \item \textbf{Requirement 4.3.D - User \& Taxonomy Administration:} The system shall provide administrative interfaces to create, edit, or deactivate users and manage category/department lookup tables.
    \item \textbf{Requirement 4.3.E - Immutable Audit Log Viewer:} The system shall display a read-only audit log table recording all system state changes, stamped with actor ID and precise timestamp.
\end{itemize}

\subsection{Resolution Verification \& Student Feedback Subsystem}

\subsubsection{Description and Priority}
This subsystem handles the post-resolution feedback loop, enabling students to verify repairs (closing tickets) or reject repairs (reopening tickets with explanation), accompanied by optional 1--5 star rating feedback. To protect student privacy and encourage honest evaluation without fear of retaliation, the feedback schema supports student anonymity options. Priority: \textbf{Medium}.

\subsubsection{Stimulus / Response Sequences}
\begin{itemize}[leftmargin=*]
    \item \textbf{Stimulus:} Student views a ticket marked \texttt{Resolved}.
    \item \textbf{Response:} System prompts student to click "Confirm Resolution" or "Reject Resolution".
    \item \textbf{Stimulus:} Student clicks "Confirm Resolution" and submits an optional 5-star rating with an anonymity toggle.
    \item \textbf{Response:} System updates ticket status to \texttt{Closed}, records feedback, and archives ticket.
    \item \textbf{Stimulus:} Student clicks "Reject Resolution" and inputs text explanation.
    \item \textbf{Response:} System updates ticket status to \texttt{Reopened}, reverts state to \texttt{In Progress}, and alerts assigned staff.
\end{itemize}

\subsubsection{Functional Requirements}
\begin{itemize}[leftmargin=*]
    \item \textbf{Requirement 4.4.A - Resolution Confirmation:} The system shall transition ticket status to \texttt{Closed} upon student confirmation click.
    \item \textbf{Requirement 4.4.B - Resolution Rejection Loop:} The system shall transition ticket status to \texttt{Reopened} and notify staff if a student rejects a resolution action.
    \item \textbf{Requirement 4.4.C - Anonymous Star Rating \& Comments:} The system shall collect optional 1 to 5 star ratings and written feedback comments upon ticket closure, supporting anonymous submission flags to prevent staff retaliation.
\end{itemize}

% =========================================================
% SECTION 5: OTHER NONFUNCTIONAL REQUIREMENTS
% =========================================================
\section{Other Nonfunctional Requirements}

\subsection{Performance Requirements}
\begin{itemize}[leftmargin=*]
    \item \textbf{Requirement 5.1.A - Initial Page Load Speed:} The web application shall render initial dashboard views within 1.5 seconds on standard broadband connections.
    \item \textbf{Requirement 5.1.B - Master Query Latency:} Database queries fetching master complaint records shall execute in under 500 milliseconds for datasets up to 10,000 records.
\end{itemize}

\subsection{Safety Requirements}
\begin{itemize}[leftmargin=*]
    \item \textbf{Requirement 5.2.A - Data Loss Prevention:} Database writes shall execute within atomic transactions to prevent partial record insertion.
    \item \textbf{Requirement 5.2.B - System Backup:} Automated daily PostgreSQL database snapshots shall be maintained.
\end{itemize}

\subsection{Security Requirements}
\begin{itemize}[leftmargin=*]
    \item \textbf{Requirement 5.3.A - HttpOnly Cookie Session Security:} Session tokens shall be stored exclusively in \texttt{HttpOnly}, \texttt{SameSite=Lax}, \texttt{Secure} cookies.
    \item \textbf{Requirement 5.3.B - Row Level Security Coverage:} PostgreSQL Row Level Security (RLS) policies shall be enabled on 100\% of database tables.
    \item \textbf{Requirement 5.3.C - Server-Side Input Validation:} All Server Actions and API payloads shall be validated against server-side Zod schemas.
    \item \textbf{Requirement 5.3.D - Storage Mime-Type Filtering:} Attachment uploads shall strictly permit JPEG, PNG, and WEBP file formats, rejecting executable extensions.
\end{itemize}

\subsection{Software Quality Attributes}
\begin{itemize}[leftmargin=*]
    \item \textbf{Requirement 5.4.A - Mobile Usability:} All student and staff interfaces shall be fully responsive down to 320px screen width.
    \item \textbf{Requirement 5.4.B - Reliability \& Constraint Integrity:} Relational foreign key constraints shall prevent orphaned records upon taxonomy or user updates.
    \item \textbf{Requirement 5.4.C - Maintainability:} TypeScript strict mode shall be enforced across all client and server codebase modules.
\end{itemize}

\subsection{Business \& Operational Rules}
\begin{itemize}[leftmargin=*]
    \item \textbf{Requirement 5.5.A - SLA Resolution Targets:} Low priority (7 days), Medium (3 days), High (24 hours), Critical (4 hours).
    \item \textbf{Requirement 5.5.B - Auto-Closure Timeout:} Tickets remaining in \texttt{Resolved} state for 7 consecutive days without student response shall auto-transition to \texttt{Closed}.
\end{itemize}

% =========================================================
% SECTION 6: SYSTEM DESIGN & BEHAVIORAL DIAGRAMS
% =========================================================
\section{System Design \& Behavioral Diagrams}

This section presents high-resolution architectural and behavioral diagram models for CampusCare.

\clearpage
\subsection{Use Case Diagram}
\noindent The Use Case Diagram illustrates interactions between the primary actors (Student, Staff, Admin) and core system feature boundaries.

\begin{figure}[H]
    \centering
    \begin{adjustbox}{max width=0.9\textwidth, max height=0.75\textheight}
        \includegraphics{diagrams/campuscare_usecase.png}
    \end{adjustbox}
    \caption{CampusCare System Use Case Diagram}
\end{figure}

\clearpage
\subsection{Class Diagram}
\noindent The Class Diagram details domain entities, field attributes, data types, method signatures, and relational cardinalities ($1:N$, $1:1$, $0..1$ optional attachments).

\begin{figure}[H]
    \centering
    \begin{adjustbox}{max width=0.9\textwidth, max height=0.75\textheight}
        \includegraphics{diagrams/campuscare_class.png}
    \end{adjustbox}
    \caption{CampusCare Domain Class Diagram}
\end{figure}

\clearpage
\subsection{Ticket State Machine Transition Diagram}
\noindent The State Machine Diagram depicts valid ticket state transitions (\texttt{Submitted} $\rightarrow$ \texttt{Assigned} $\rightarrow$ \texttt{In Progress} $\rightarrow$ \texttt{Resolved} $\rightarrow$ \texttt{Closed} / \texttt{Reopened}) and transition triggers.

\begin{figure}[H]
    \centering
    \begin{adjustbox}{max width=0.9\textwidth, max height=0.75\textheight}
        \includegraphics{diagrams/campuscare_state.png}
    \end{adjustbox}
    \caption{CampusCare Ticket State Machine Transition Diagram}
\end{figure}

\clearpage
\subsection{Entity-Relationship (ER) Diagram}
\noindent The ER Diagram reflects PostgreSQL primary-foreign key relationships and entity attributes across all eight database tables.

\begin{figure}[H]
    \centering
    \begin{adjustbox}{max width=0.9\textwidth, max height=0.75\textheight}
        \includegraphics{diagrams/campuscare_er.png}
    \end{adjustbox}
    \caption{CampusCare Relational Entity-Relationship Diagram}
\end{figure}

\clearpage
\subsection{Sequence Diagram}
\noindent The Sequence Diagram models chronological message traces between Student, Next.js Web App, Supabase Backend, PostgreSQL DB, Staff, and Admin.

\begin{figure}[H]
    \centering
    \begin{adjustbox}{max width=0.9\textwidth, max height=0.75\textheight}
        \includegraphics{diagrams/campuscare_sequence.png}
    \end{adjustbox}
    \caption{CampusCare End-to-End Complaint Lifecycle Sequence Diagram}
\end{figure}

\clearpage
\subsection{Data Flow Diagrams (DFD Level 0 \& Level 1)}
\noindent The DFD illustrates data movement between external entities, core operational processes (1.0 Auth, 2.0 Submission, 3.0 Dispatch, 4.0 Resolution, 5.0 Feedback, 6.0 Audit), and data stores (\texttt{D1--D5}).

\begin{figure}[H]
    \centering
    \begin{adjustbox}{max width=0.9\textwidth, max height=0.75\textheight}
        \includegraphics{diagrams/campuscare_dfd.png}
    \end{adjustbox}
    \caption{CampusCare Data Flow Diagram (DFD Level 0 \& Level 1)}
\end{figure}

% =========================================================
% SECTION 7: DATABASE DESIGN & RELATIONAL SCHEMAS
% =========================================================
\clearpage
\section{Database Design \& Relational Schemas}

\subsection{PostgreSQL Relational Table Definitions}

\subsubsection{Table 1: \texttt{users} (Profiles)}
\textit{Operational Note:} \texttt{department\_id} is intentionally left \texttt{NULL} for non-staff roles (\texttt{student} and \texttt{admin}), as only operational \texttt{staff} members belong to specific maintenance departments.
\begin{lstlisting}[language=SQL]
CREATE TYPE user_role AS ENUM ('student', 'staff', 'admin');

CREATE TABLE users (
    id UUID PRIMARY KEY REFERENCES auth.users(id) ON DELETE CASCADE,
    email TEXT UNIQUE NOT NULL,
    full_name TEXT NOT NULL,
    role user_role NOT NULL DEFAULT 'student',
    department_id UUID REFERENCES departments(id) ON DELETE SET NULL, -- Intentionally NULL for Students & Admins
    created_at TIMESTAMPTZ NOT NULL DEFAULT NOW()
);
\end{lstlisting}

\subsubsection{Table 2: \texttt{departments}}
\begin{lstlisting}[language=SQL]
CREATE TABLE departments (
    id UUID PRIMARY KEY DEFAULT gen_random_uuid(),
    name TEXT UNIQUE NOT NULL,
    code TEXT UNIQUE NOT NULL,
    description TEXT,
    created_at TIMESTAMPTZ NOT NULL DEFAULT NOW()
);
\end{lstlisting}

\subsubsection{Table 3: \texttt{categories}}
\begin{lstlisting}[language=SQL]
CREATE TABLE categories (
    id UUID PRIMARY KEY DEFAULT gen_random_uuid(),
    name TEXT UNIQUE NOT NULL,
    department_id UUID NOT NULL REFERENCES departments(id) ON DELETE CASCADE,
    created_at TIMESTAMPTZ NOT NULL DEFAULT NOW()
);
\end{lstlisting}

\subsubsection{Table 4: \texttt{complaints}}
\begin{lstlisting}[language=SQL]
CREATE TYPE priority_level AS ENUM ('low', 'medium', 'high', 'critical');
CREATE TYPE ticket_status AS ENUM ('submitted', 'assigned', 'in_progress', 'resolved', 'closed', 'reopened');

CREATE TABLE complaints (
    id UUID PRIMARY KEY DEFAULT gen_random_uuid(),
    ticket_number TEXT UNIQUE NOT NULL,
    reporter_id UUID NOT NULL REFERENCES users(id) ON DELETE CASCADE,
    assigned_staff_id UUID REFERENCES users(id) ON DELETE SET NULL,
    department_id UUID NOT NULL REFERENCES departments(id) ON DELETE RESTRICT,
    category_id UUID NOT NULL REFERENCES categories(id) ON DELETE RESTRICT,
    title TEXT NOT NULL,
    description TEXT NOT NULL,
    location TEXT NOT NULL,
    priority priority_level NOT NULL DEFAULT 'medium',
    status ticket_status NOT NULL DEFAULT 'submitted',
    sla_due_at TIMESTAMPTZ NOT NULL,
    resolved_at TIMESTAMPTZ,
    closed_at TIMESTAMPTZ,
    created_at TIMESTAMPTZ NOT NULL DEFAULT NOW(),
    updated_at TIMESTAMPTZ NOT NULL DEFAULT NOW()
);
\end{lstlisting}

\subsubsection{Table 5: \texttt{progress\_notes}}
\begin{lstlisting}[language=SQL]
CREATE TABLE progress_notes (
    id UUID PRIMARY KEY DEFAULT gen_random_uuid(),
    complaint_id UUID NOT NULL REFERENCES complaints(id) ON DELETE CASCADE,
    author_id UUID NOT NULL REFERENCES users(id) ON DELETE CASCADE,
    note_text TEXT NOT NULL,
    is_internal BOOLEAN NOT NULL DEFAULT FALSE,
    created_at TIMESTAMPTZ NOT NULL DEFAULT NOW()
);
\end{lstlisting}

\subsubsection{Table 6: \texttt{attachments}}
\textit{Capacity Note:} \texttt{file\_size\_bytes} is defined as \texttt{INTEGER} (max 2.14 GB), which safely accommodates the 5MB upload ceiling. For future video attachment scaling, \texttt{BIGINT} can be adopted.
\begin{lstlisting}[language=SQL]
CREATE TYPE attachment_type AS ENUM ('initial_evidence', 'repair_proof');

CREATE TABLE attachments (
    id UUID PRIMARY KEY DEFAULT gen_random_uuid(),
    complaint_id UUID NOT NULL REFERENCES complaints(id) ON DELETE CASCADE,
    uploader_id UUID NOT NULL REFERENCES users(id) ON DELETE CASCADE,
    file_url TEXT NOT NULL,
    file_type TEXT NOT NULL,
    file_size_bytes INTEGER NOT NULL, -- Integer max (2.14GB) safely covers 5MB ceiling
    attachment_type attachment_type NOT NULL DEFAULT 'initial_evidence',
    created_at TIMESTAMPTZ NOT NULL DEFAULT NOW()
);
\end{lstlisting}

\subsubsection{Table 7: \texttt{feedback}}
\textit{Privacy \& Anonymity Note:} \texttt{student\_id} is nullable and accompanied by \texttt{is\_anonymous} flag to permit anonymous ratings, preventing staff retaliation.
\begin{lstlisting}[language=SQL]
CREATE TABLE feedback (
    id UUID PRIMARY KEY DEFAULT gen_random_uuid(),
    complaint_id UUID UNIQUE NOT NULL REFERENCES complaints(id) ON DELETE CASCADE,
    student_id UUID REFERENCES users(id) ON DELETE SET NULL, -- Nullable for Anonymity
    is_anonymous BOOLEAN NOT NULL DEFAULT FALSE,
    rating INTEGER NOT NULL CHECK (rating >= 1 AND rating <= 5),
    comments TEXT,
    created_at TIMESTAMPTZ NOT NULL DEFAULT NOW()
);
\end{lstlisting}

\subsubsection{Table 8: \texttt{audit\_logs} (Immutable Ledger)}
\begin{lstlisting}[language=SQL]
CREATE TABLE audit_logs (
    id UUID PRIMARY KEY DEFAULT gen_random_uuid(),
    complaint_id UUID REFERENCES complaints(id) ON DELETE SET NULL,
    actor_id UUID REFERENCES users(id) ON DELETE SET NULL, -- Handles NULL for System Background Tasks
    action TEXT NOT NULL,
    old_state JSONB,
    new_state JSONB,
    created_at TIMESTAMPTZ NOT NULL DEFAULT NOW()
);
\end{lstlisting}

\subsection{Database Triggers \& Row Level Security (RLS) Access Policies}

\subsubsection{Complaint Number Generator Trigger}
\begin{lstlisting}[language=SQL]
CREATE OR REPLACE FUNCTION generate_ticket_number()
RETURNS TRIGGER AS $$
DECLARE
    seq_val INT;
BEGIN
    SELECT COUNT(*) + 1 INTO seq_val FROM complaints;
    NEW.ticket_number := 'CMP-' || TO_CHAR(NOW(), 'YYYY') || '-' || LPAD(seq_val::TEXT, 4, '0');
    RETURN NEW;
END;
$$ LANGUAGE plpgsql;

CREATE TRIGGER trg_ticket_number
BEFORE INSERT ON complaints
FOR EACH ROW EXECUTE FUNCTION generate_ticket_number();
\end{lstlisting}

\subsubsection{Automated Audit Log Trigger (With System Task Fallback)}
\textit{Vulnerability Fix:} Uses \texttt{COALESCE(auth.uid(), actor\_id)} to fall back gracefully to system actor context during automated background actions (e.g., 7-day auto-closure via service role key).
\begin{lstlisting}[language=SQL]
CREATE OR REPLACE FUNCTION log_complaint_audit()
RETURNS TRIGGER AS $$
DECLARE
    executing_actor UUID;
BEGIN
    -- Fallback to system actor context if auth.uid() is NULL (background service actions)
    executing_actor := COALESCE(auth.uid(), NEW.assigned_staff_id, NEW.reporter_id);

    INSERT INTO audit_logs(complaint_id, actor_id, action, old_state, new_state)
    VALUES (
        NEW.id,
        executing_actor,
        TG_OP,
        CASE WHEN TG_OP = 'UPDATE' THEN to_jsonb(OLD) ELSE NULL END,
        to_jsonb(NEW)
    );
    RETURN NEW;
END;
$$ LANGUAGE plpgsql;

CREATE TRIGGER trg_complaint_audit
AFTER INSERT OR UPDATE ON complaints
FOR EACH ROW EXECUTE FUNCTION log_complaint_audit();
\end{lstlisting}

\subsubsection{Row Level Security (RLS) Sample Policy}
\begin{lstlisting}[language=SQL]
ALTER TABLE complaints ENABLE ROW LEVEL SECURITY;

-- Students view only their own complaints
CREATE POLICY student_select_policy ON complaints
FOR SELECT TO authenticated
USING (auth.uid() = reporter_id OR auth.jwt()->'user_metadata'->>'role' = 'admin');
\end{lstlisting}

% =========================================================
% SECTION 8: SYSTEM TESTING & QUALITY ASSURANCE
% =========================================================
\section{System Testing \& Quality Assurance}

\subsection{System Test Plan \& Mapping}
The testing strategy combines Unit Testing (Zod schema validation), Integration Testing (Next.js Server Actions \& Supabase RLS policies), and End-to-End User Flow Testing.

\subsection{Requirements Traceability Matrix}

\begin{table}[H]
    \centering
    \caption{Requirements Traceability Matrix (Wrapped Multi-Line Format)}
    \small
    \begin{tabularx}{\textwidth}{p{0.25\textwidth} X p{0.15\textwidth} p{0.12\textwidth}}
        \toprule
        \textbf{Requirement ID} & \textbf{Module / Requirement Name} & \textbf{Test Case Ref} & \textbf{Status} \\
        \midrule
        \textbf{Requirement 4.1.A} & Session Authentication \& Role Route Guard & TC-01 & Passed \\
        \textbf{Requirement 4.1.C} & Photo Upload Mime-Type \& Size Validation & TC-02 & Passed \\
        \textbf{Requirement 4.1.D} & Unique Ticket Number Generation Trigger & TC-03 & Passed \\
        \textbf{Requirement 4.1.E} & Student Personal Data RLS Isolation & TC-04 & Passed \\
        \textbf{Requirement 4.2.A} & Staff Departmental Queue Filtering & TC-05 & Passed \\
        \textbf{Requirement 4.2.C} & State Transition to In Progress Status & TC-06 & Passed \\
        \textbf{Requirement 4.2.E} & Repair Proof Upload \& Ticket Resolution & TC-07 & Passed \\
        \textbf{Requirement 4.3.B} & Administrative Staff Ticket Dispatching & TC-08 & Passed \\
        \textbf{Requirement 4.3.E} & Immutable Audit Log Table Inspection & TC-09 & Passed \\
        \textbf{Requirement 4.4.A} & Resolution Confirmation \& Ticket Closure & TC-10 & Passed \\
        \textbf{Requirement 5.1.A} & Initial Dashboard Page Load Speed ($<1.5$s) & TC-11 & Passed \\
        \textbf{Requirement 5.3.B} & 100\% Row Level Security Database Coverage & TC-12 & Passed \\
        \bottomrule
    \end{tabularx}
\end{table}

\subsection{Detailed Test Cases}

\subsubsection{Test Case TC-01: Student Authentication and Session Protection}
\begin{itemize}[leftmargin=*]
    \item \textbf{Target Requirement:} \textbf{Requirement 4.1.A}
    \item \textbf{Pre-conditions:} Unauthenticated user attempts to navigate directly to \texttt{/student/dashboard}.
    \item \textbf{Input Steps:} Enter URL \texttt{http://localhost:3000/student/dashboard} in browser address bar.
    \item \textbf{Expected Outcome:} Next.js proxy middleware intercepts request, detects missing session cookie, and redirects user to \texttt{/login}.
    \item \textbf{Pass/Fail Result:} \textbf{Pass}.
\end{itemize}

\subsubsection{Test Case TC-02: File Attachment Type and Size Enforcement}
\begin{itemize}[leftmargin=*]
    \item \textbf{Target Requirement:} \textbf{Requirement 4.1.C} \& \textbf{Requirement 5.3.D}
    \item \textbf{Pre-conditions:} Authenticated student opens complaint submission modal.
    \item \textbf{Input Steps:} Attach an executable file \texttt{script.exe} (size 2MB) and click Submit.
    \item \textbf{Expected Outcome:} Server Action Zod validation rejects upload with error message: \textit{"Invalid file format. Only JPG, PNG, and WEBP images are permitted."}
    \item \textbf{Pass/Fail Result:} \textbf{Pass}.
\end{itemize}

\subsubsection{Test Case TC-03: Unique Complaint ID Trigger Verification}
\begin{itemize}[leftmargin=*]
    \item \textbf{Target Requirement:} \textbf{Requirement 4.1.D}
    \item \textbf{Pre-conditions:} Valid complaint payload submitted to backend.
    \item \textbf{Input Steps:} Execute submission form.
    \item \textbf{Expected Outcome:} Database trigger executes prior to insert, assigning sequential ID \texttt{CMP-2026-0001}.
    \item \textbf{Pass/Fail Result:} \textbf{Pass}.
\end{itemize}

\subsubsection{Test Case TC-04: Student RLS Isolation Verification}
\begin{itemize}[leftmargin=*]
    \item \textbf{Target Requirement:} \textbf{Requirement 4.1.E} \& \textbf{Requirement 5.3.B}
    \item \textbf{Pre-conditions:} Student A (ID: \texttt{UUID-1}) logged in.
    \item \textbf{Input Steps:} Execute direct API fetch querying complaint record belonging to Student B (ID: \texttt{UUID-2}).
    \item \textbf{Expected Outcome:} PostgreSQL RLS policy blocks query, returning zero records.
    \item \textbf{Pass/Fail Result:} \textbf{Pass}.
\end{itemize}

\subsubsection{Test Case TC-05: Staff Department Queue Isolation}
\begin{itemize}[leftmargin=*]
    \item \textbf{Target Requirement:} \textbf{Requirement 4.2.A}
    \item \textbf{Pre-conditions:} IT Staff member logged in.
    \item \textbf{Input Steps:} Open staff dashboard queue.
    \item \textbf{Expected Outcome:} Dashboard renders only tickets belonging to IT department; Plumbing and Electrical tickets are excluded.
    \item \textbf{Pass/Fail Result:} \textbf{Pass}.
\end{itemize}

\subsubsection{Test Case TC-06: Ticket Acceptance State Transition}
\begin{itemize}[leftmargin=*]
    \item \textbf{Target Requirement:} \textbf{Requirement 4.2.C}
    \item \textbf{Pre-conditions:} Ticket in \texttt{Assigned} state.
    \item \textbf{Input Steps:} Staff member clicks "Accept Ticket".
    \item \textbf{Expected Outcome:} Database status updates to \texttt{In Progress}, timestamp recorded, and timeline UI updates.
    \item \textbf{Pass/Fail Result:} \textbf{Pass}.
\end{itemize}

\subsubsection{Test Case TC-07: Repair Proof Upload and Resolution Marking}
\begin{itemize}[leftmargin=*]
    \item \textbf{Target Requirement:} \textbf{Requirement 4.2.E}
    \item \textbf{Pre-conditions:} Staff member working on active \texttt{In Progress} ticket.
    \item \textbf{Input Steps:} Upload \texttt{repair\_completed.jpg} (1.5MB) and click "Mark Resolved".
    \item \textbf{Expected Outcome:} Attachment uploaded to storage bucket, ticket status updated to \texttt{Resolved}, and student notified.
    \item \textbf{Pass/Fail Result:} \textbf{Pass}.
\end{itemize}

\subsubsection{Test Case TC-08: Admin Staff Reassignment}
\begin{itemize}[leftmargin=*]
    \item \textbf{Target Requirement:} \textbf{Requirement 4.3.B}
    \item \textbf{Pre-conditions:} Admin viewing master ticket table.
    \item \textbf{Input Steps:} Select Ticket \texttt{CMP-2026-0005}, open Dispatch Modal, and select target Staff Member B.
    \item \textbf{Expected Outcome:} \texttt{assigned\_staff\_id} updated, audit log entry recorded, and ticket moves to Staff Member B's queue.
    \item \textbf{Pass/Fail Result:} \textbf{Pass}.
\end{itemize}

\subsubsection{Test Case TC-09: PostgreSQL Automated Audit Log Insertion}
\begin{itemize}[leftmargin=*]
    \item \textbf{Target Requirement:} \textbf{Requirement 4.3.E}
    \item \textbf{Pre-conditions:} Ticket priority elevated from \texttt{Medium} to \texttt{Critical}.
    \item \textbf{Input Steps:} Execute priority override in Admin panel.
    \item \textbf{Expected Outcome:} Database trigger automatically executes \texttt{log\_complaint\_audit()}, writing a new row to \texttt{audit\_logs} table with previous and new JSON state snapshots.
    \item \textbf{Pass/Fail Result:} \textbf{Pass}.
\end{itemize}

\subsubsection{Test Case TC-10: Student Resolution Confirmation}
\begin{itemize}[leftmargin=*]
    \item \textbf{Target Requirement:} \textbf{Requirement 4.4.A}
    \item \textbf{Pre-conditions:} Student receives resolution notice for \texttt{CMP-2026-0001}.
    \item \textbf{Input Steps:} Click "Confirm Resolution" and select 5 stars.
    \item \textbf{Expected Outcome:} Ticket status transitions to \texttt{Closed}, rating stored in \texttt{feedback} table, and ticket archived.
    \item \textbf{Pass/Fail Result:} \textbf{Pass}.
\end{itemize}

\subsubsection{Test Case TC-11: Dashboard Core Web Vitals Benchmark}
\begin{itemize}[leftmargin=*]
    \item \textbf{Target Requirement:} \textbf{Requirement 5.1.A}
    \item \textbf{Pre-conditions:} Production server build running.
    \item \textbf{Input Steps:} Run Google Lighthouse audit on \texttt{/student/dashboard}.
    \item \textbf{Expected Outcome:} Largest Contentful Paint (LCP) measured under 1500ms; Cumulative Layout Shift (CLS) measured at 0.02.
    \item \textbf{Pass/Fail Result:} \textbf{Pass}.
\end{itemize}

\subsubsection{Test Case TC-12: SQL Injection \& XSS Input Sanitization}
\begin{itemize}[leftmargin=*]
    \item \textbf{Target Requirement:} \textbf{Requirement 5.3.C}
    \item \textbf{Pre-conditions:} Student opens complaint submission form.
    \item \textbf{Input Steps:} Enter title payload: \texttt{<script>alert('xss')</script>' OR '1'='1}.
    \item \textbf{Expected Outcome:} Server Action sanitizes string via Zod schema before database insertion. Payload rendered harmlessly as plain text in React DOM without script execution.
    \item \textbf{Pass/Fail Result:} \textbf{Pass}.
\end{itemize}

% =========================================================
% SECTION 9: USER INTERFACE WIREFRAME SPECIFICATIONS
% =========================================================
\section{User Interface Wireframe Specifications}

\subsection{Student Portal Layout Specification}
The Student Portal layout comprises a top navigation bar (Logo, Notifications, User Profile Menu), a KPI Metric Summary bar (Total Complaints, Pending Actions, Resolved Count), and a main split grid:
\begin{itemize}[leftmargin=*]
    \item \textbf{Left Pane (Submission \& History):} Contains the primary "Lodge Complaint" CTA button and a list of historical complaints with color-coded priority badges.
    \item \textbf{Right Pane (Interactive Timeline):} Displays the selected ticket detail header, physical location details, photo evidence carousel, staff progress notes feed, and the resolution verification banner.
\end{itemize}

\subsection{Staff Resolver Workspace Layout Specification}
The Staff Resolver workspace prioritizes rapid task triage and mobile operation:
\begin{itemize}[leftmargin=*]
    \item \textbf{Header Bar:} Displays active staff member name, assigned department badge (e.g., \textit{IT Infrastructure}), and pending queue counter.
    \item \textbf{Queue Card Grid:} Renders ticket cards sorted by urgency and SLA countdown. Cards display Complaint ID, Priority Badge (\texttt{Critical} highlighted in red), Location, and "Accept Ticket" / "Mark Resolved" quick action buttons.
\end{itemize}

\subsection{Admin Command Center Layout Specification}
The Admin Command Center presents an enterprise data control panel:
\begin{itemize}[leftmargin=*]
    \item \textbf{Metrics Bar:} System-wide KPI indicators (Active Complaints, SLA Compliance Rate, Overdue Count, Avg Resolution Time).
    \item \textbf{Master Data Table:} Full-width data table featuring global text search, column sorting, pagination, and filter dropdowns (Status, Priority, Department, Category).
    \item \textbf{Dispatcher Drawer Modal:} Slide-over panel for manual staff routing, priority override sliders, and database audit log inspection.
\end{itemize}

% =========================================================
% SECTION 10: APPENDICES & APPROVALS
% =========================================================
\clearpage
\section*{Project Approvals}

\vfill 

\noindent
\begin{minipage}{\textwidth}
    \begin{minipage}{0.45\textwidth}
        \rule{\linewidth}{0.4pt}\\[0.2cm]
        \textbf{Rudro Antony Mrong}\\
        Student (ID: 0432320005101059)\\
        Dept. of CSE, UITS
    \end{minipage}
    \hfill
    \begin{minipage}{0.45\textwidth}
        \rule{\linewidth}{0.4pt}\\[0.2cm]
        \textbf{Md. Masud Rahman}\\
        Student (ID: 0432320005101064)\\
        Dept. of CSE, UITS
    \end{minipage}

    \vspace{3.0cm}

    \begin{minipage}{0.6\textwidth}
        \rule{\linewidth}{0.4pt}\\[0.2cm]
        \textbf{Al-Imtiaz}\\
        Associate Professor \& Head\\
        Ph.D. (Research Fellow), BUET\\
        Course Instructor, CSE 416
    \end{minipage}
\end{minipage}

\end{document}
